% Build with latexmk from this directory. The PDF is written to build/main.pdf.
\documentclass[a4paper,12pt]{article}

% =========================
% Layout and paragraph style
% =========================
\usepackage[left=3.18cm,right=3.18cm,top=2.54cm,bottom=2.54cm]{geometry}
\usepackage{setspace}
\setstretch{1.25} % Line spacing. Common choices: 1.2 / 1.25 / 1.5.

\usepackage{parskip}        % Adds vertical space between paragraphs.
\setlength{\parindent}{0em} % Set to non-zero if you prefer first-line indentation.

% =========================
% Font setup
% =========================
% Use XeLaTeX so English and Chinese fonts can be configured directly.
\usepackage{fontspec}
\usepackage{xeCJK}

% Use the original report fonts when installed, with TeX Live alternatives.
\IfFontExistsTF{Times New Roman}
  {\setmainfont{Times New Roman}}
  {\setmainfont{texgyretermes}[Extension=.otf,
     UprightFont=*-regular, BoldFont=*-bold,
     ItalicFont=*-italic, BoldItalicFont=*-bolditalic]}

\ExplSyntaxOn
\newcommand*{\ReportCJKFont}{\sys_if_platform_windows:TF {SimSun} {Songti~SC}}
\ExplSyntaxOff
\IfFontExistsTF{\ReportCJKFont}
  {\setCJKmainfont{\ReportCJKFont}[Script=Default]}
  {\setCJKmainfont{FandolSong-Regular.otf}[Script=Default,
     BoldFont=FandolSong-Bold.otf, ItalicFont=FandolKai-Regular.otf]}

% Reuse the configured CJK font in sans and monospaced labels.
\renewcommand*{\CJKsfdefault}{\CJKrmdefault}
\renewcommand*{\CJKttdefault}{\CJKrmdefault}

% =========================
% Common packages
% =========================
\usepackage{graphicx}
\usepackage{xcolor}
\usepackage{eso-pic}
\usepackage{etoolbox}
\usepackage{tabularray}
\usepackage{amsmath,amssymb}
\usepackage{booktabs}
\usepackage{hyperref}
\usepackage{pdfpages}
\usepackage[backend=biber,style=numeric,sorting=none]{biblatex}
\addbibresource{references.bib}

% Default folder for \includegraphics.
\graphicspath{{figures/}}

% Keep hyperlinks visually simple in academic documents.
\hypersetup{
  colorlinks=true,
  linkcolor=black,
  urlcolor=black,
  citecolor=black
}

\sloppy % Allow LaTeX to adjust spacing to prevent overfull lines.

% Copy the example into local-materials/course-settings.tex to personalize it.
\IfFileExists{local-materials/course-settings.tex}
  {\input{local-materials/course-settings}}
  {% Copy this file to local-materials/course-settings.tex for a real assignment.
% The local copy replaces this example. Long text wraps at spaces in the table.
\newcommand{\ReportInstitution}{EXAMPLE UNIVERSITY}
\newcommand{\CoverCourseCode}{DEMO101}
\newcommand{\CoverCourseName}{Example Course}
\newcommand{\CoverLecturer}{Example Lecturer}
\newcommand{\CoverSession}{Example Term}
\newcommand{\CoverAssessment}{Example Assignment Title}
\newcommand{\CoverDueDate}{DD / MM / YYYY}
\newcommand{\ReportAuthors}{Example Student} % Title block in the report body.

% Add one row per student. Remove unused rows.
\newcommand{\CoverStudents}{%
  \CoverStudent{ABC1234567}{Example Student}%
}

% Supply one signature per student, in the same order. The example draws a label.
% In the local copy, replace the placeholder line with one or more image lines.
% PNG, JPEG and PDF are supported. The optional height defaults to 1.7 cm.
\newcommand{\CoverSignatures}{%
  \CoverSignaturePlaceholder
  % \CoverSignatureImage{local-materials/signatures/student-one.png}%
  % \CoverSignatureImage[2cm]{local-materials/signatures/student-two.png}%
}

% The logo keeps its aspect ratio in the cover's square image slot.
\newcommand{\CoverLogoFile}{local-materials/logo.pdf}
% Supply the course's declaration as paragraphs in local-materials/declaration.tex.
\newcommand{\CoverDeclarationTitle}{Course Declaration}
\newcommand{\CoverDeclarationDate}{\today}

% Leave these fields empty for the marker.
\newcommand{\CoverFeedback}{}
\newcommand{\CoverMark}{}
}

% =========================
% Header and footer
% =========================
\usepackage{fancyhdr}
\pagestyle{fancy}
\fancyhf{} % Clear all default header/footer fields first.

\fancyhead[C]{\fontsize{12pt}{14pt}\selectfont \ReportInstitution}
\fancyfoot[C]{\fontsize{12pt}{14pt}\selectfont \thepage}

\renewcommand{\headrulewidth}{0pt} % Change to 0.4pt if you want a header rule.
\renewcommand{\footrulewidth}{0pt}
\setlength{\headheight}{15pt}      % Prevents fancyhdr headheight warnings.

% =========================
% Table of contents depth
% =========================
\setcounter{tocdepth}{2} % 1=section, 2=subsection, 3=subsubsection

\begin{document}

% !TeX root = ../main.tex
% Change course data in a local copy of course-settings-example.tex.
\begingroup

\newcommand{\CoverStudentRows}{}
\newcommand{\CoverStudent}[2]{%
  \appto\CoverStudentRows{& & {#1} & {#2} \\}}
\CoverStudents
\ifdefempty{\CoverStudentRows}{\CoverStudent{}{}}{}

\newcommand{\CoverSignatureRows}{}
\newcommand{\CoverSignatureImage}[2][1.7cm]{%
  \appto\CoverSignatureRows{%
    \hbox{\includegraphics[width=350bp,height=#1,keepaspectratio]{#2}}%
    \vskip 5bp}}
\newcommand{\CoverSignaturePlaceholder}{%
  \appto\CoverSignatureRows{%
    \hbox{\begin{picture}(128,48)
      \linethickness{0.48bp}%
      \put(0.24,0.24){\framebox(127.52,47.52){\normalfont\sffamily\small SIGNATURE}}%
    \end{picture}}%
    \vskip 5bp}}
\CoverSignatures

% PostScript points, measured from the upper left of the A4 page.
\newcommand{\CoverPlace}[3]{\put(#1,-#2){#3}}
\newcommand{\CoverPage}[1]{%
  \thispagestyle{empty}%
  \AddToShipoutPictureFG*{\AtPageUpperLeft{%
    \setlength{\unitlength}{1bp}%
    \renewcommand{\baselinestretch}{1}%
    \normalfont\fontsize{12bp}{14bp}\selectfont
    \CoverPlace{297.05}{47.16}{\makebox[0pt][c]{\ReportInstitution}}%
    #1}}%
  \null\clearpage}

% Use a simple typographic placeholder until a local course logo is supplied.
\newcommand{\CoverLogo}{%
  \begin{picture}(108.75,107.95)
    \IfFileExists{\CoverLogoFile}{%
      \put(0,0){\makebox(108.75,107.95){%
        \includegraphics[width=108.75bp,height=107.95bp,keepaspectratio]{\CoverLogoFile}}}%
    }{%
      \linethickness{0.48bp}%
      \put(0.24,0.24){\framebox(108.27,107.47){\normalfont\sffamily\large LOGO}}%
    }%
  \end{picture}}

\newcommand{\CoverFeedbackBox}{%
  \begin{picture}(415.52,161.55)
    \linethickness{0.48bp}%
    \put(0.24,0.24){\framebox(415.04,161.07){}}%
    \put(5.40,148.23){Feedback from Lecturer:}%
    \put(5.40,129.75){\parbox[t]{326bp}{\raggedright\CoverFeedback}}%
    % The outer frame supplies the right and bottom of the 72 bp mark box.
    \put(343.52,0){\line(0,1){72}}%
    \put(343.28,72){\line(1,0){72.24}}%
    \put(350.45,56.07){Mark:}%
    \put(350.45,36.90){\CoverMark}%
  \end{picture}}

\newcommand{\CoverForm}{%
  \vbox{\offinterlineskip
    \hbox{%
      \begin{tblr}[expand=\CoverStudentRows]{
        colspec={Q[l,wd=124.70bp] Q[c,wd=6.00bp] Q[l,wd=83.30bp] Q[l,wd=165.64bp]},
        colsep=5.40bp, rowsep=0pt,
        column{1}={rightsep=0.40bp}, column{2}={leftsep=3.64bp,rightsep=3.64bp},
        column{3}={leftsep=5.18bp,rightsep=5.62bp},
        rows={m,ht=18.00bp}, row{1-6}={ht=21.62bp,abovesep=1.85bp,belowsep=-1.85bp},
        row{8-Y}={abovesep=1.90bp,belowsep=-1.90bp}, row{Z}={ht=18.43bp},
        % Allow the merged value cells to wrap inside the original ruled area.
        cell{1-7}{3}={c=2}{l,wd=260bp}, row{8}={halign=c}, cell{8}{1}={l},
        hline{2-Y}={3-4}{0.48bp},
        vline{3,4}={8-Y}{0.48bp}, vline{5}={8-Z}{0.48bp},
      }
        Course Code         & : & \CoverCourseCode & \\
        Course Name         & : & \CoverCourseName & \\
        Lecturer            & : & \CoverLecturer   & \\
        Academic Session    & : & \CoverSession    & \\
        Assessment Title    & : & \CoverAssessment & \\
        Submission Due Date & : & \CoverDueDate    & \\
                            &   &                  & \\
        \raisebox{1.08bp}{Prepared by} & : & Student ID & Student Name \\
        \CoverStudentRows
                            &   &                  &
      \end{tblr}}%
    \hbox{\CoverFeedbackBox}}}

% Stack signature images in student order. Keep a blank line before the date.
\newcommand{\CoverSignatureBlock}{%
  \vtop{\offinterlineskip
    \hbox{\strut Signature:}%
    \ifdefempty{\CoverSignatureRows}
      {\vskip 1.7cm}
      {\CoverSignatureRows\unskip}%
    \vskip 20.7bp
    \hbox{\strut Date:}%
    \hbox{\strut\CoverDeclarationDate}}}

\newcommand{\CoverDeclarationFile}{local-materials/declaration.tex}
\newsavebox{\CoverDeclarationText}
\newlength{\CoverSignatureOffset}
\newcommand{\CoverDeclaration}{%
  \CoverPlace{297.66}{85.08}{\makebox[0pt][c]{%
    \fontsize{14.04bp}{16bp}\selectfont\bfseries
    \IfFileExists{\CoverDeclarationFile}{\CoverDeclarationTitle}{Sample Declaration}}}%
  \sbox{\CoverDeclarationText}{\parbox[t]{415.30bp}{%
    \fontsize{12bp}{20.7bp}\selectfont
    \setlength{\parindent}{0pt}%
    \setlength{\parskip}{20.7bp}%
    \hyphenpenalty=10000
    \frenchspacing
    \IfFileExists{\CoverDeclarationFile}{\input{\CoverDeclarationFile}}{%
      Replace this paragraph with the declaration supplied for your assignment.

      Check the wording against the current course instructions before signing.}}}%
  \CoverPlace{90.02}{148.70}{\usebox{\CoverDeclarationText}}%
  % Keep the standard signature position for short declarations. Longer text
  % pushes the block down so signatures do not overlap the final paragraph.
  \setlength{\CoverSignatureOffset}{\dimexpr148.70bp+\dp\CoverDeclarationText+32bp\relax}%
  \ifdim\CoverSignatureOffset<397.13bp \setlength{\CoverSignatureOffset}{397.13bp}\fi
  \CoverPlace{90.02}{\CoverSignatureOffset}{%
    \fontsize{12bp}{20.7bp}\selectfont\CoverSignatureBlock}}

\CoverPage{%
  \CoverPlace{242.90}{179.95}{\CoverLogo}%
  \CoverPlace{90.02}{193.80}{\fontsize{14.04bp}{16bp}\selectfont
    \raisebox{-\height}{\CoverForm}}}
\CoverPage{\CoverDeclaration}
\endgroup


% Title block for the main document body.
\begin{center}
{\Large\bfseries \CoverAssessment}

\vspace{0.4em}

\ReportAuthors

\vspace{0.2em}

\today

\end{center}

% Abstract is maintained as a separate file for easier editing.
% !TeX root = ../main.tex
\section*{Abstract}

State the question your report addresses, how you investigated it, and what you found.

\textbf{Keywords:} Example topic, example method

\newpage


% Generate the table of contents after the abstract.
\tableofcontents
\newpage

% Main content sections.
% !TeX root = ../main.tex
\section{Introduction}

Introduce the assignment question and explain why it matters. Cite the sources you
use to establish the background. These two fictional sources illustrate how to
cite a journal article~\parencite{example-article} and a book~\parencite{example-book}.

\subsection{Figures and tables}

Replace the sample in Figure~\ref{fig:example} with your own illustration.
Use its caption to explain what the reader should look for.

\begin{figure}[htbp]
  \centering
  \includegraphics[width=0.4\linewidth]{example-figure.png}
  \caption{Sample illustration}
  \label{fig:example}
\end{figure}

Table~\ref{tab:parameters} gives two example inputs. Replace the labels and
values with those from your assignment, including units where appropriate.

% !TeX root = ../main.tex
\begin{table}[htbp]
  \centering
  \caption{Example input values}
  \label{tab:parameters}
  \begin{tabular}{lr}
    \toprule
    Input & Sample value \\
    \midrule
    Example input A & 10 \\
    Example input B & 20 \\
    \bottomrule
  \end{tabular}
\end{table}


% % !TeX root = ../main.tex
\section{Results}

Present your observations here. Give each reported value a unit where applicable,
and refer to the figure or table that supports it.
 % Uncomment when the file is ready to include.
% !TeX root = ../main.tex
\section{Conclusion}

Answer the question from the introduction. Refer to the result that best
supports your answer and state a limitation that affects its interpretation.


% Biber supplies the bibliography; bibintoc adds its table-of-contents entry.
\clearpage
\printbibliography[heading=bibintoc,title={References}]

% Remove this block in assignments that do not require a rubric.
\IfFileExists{local-materials/rubric.pdf}
  {\includepdf[pages=-]{local-materials/rubric.pdf}}
  {% !TeX root = ../main.tex
% This original sample is used when local-materials/rubric.pdf is absent.
\clearpage
\begingroup
\thispagestyle{empty}
\definecolor{rubricCrimson}{HTML}{A00018}
\AddToShipoutPictureBG*{%
  \AtPageCenter{\makebox(0,0){%
    \rotatebox{54}{\resizebox{1.33\paperwidth}{!}{%
      \sffamily\bfseries\color{rubricCrimson!65!white}EXAMPLE}}}}}

\setlength{\parindent}{0pt}
\setlength{\parskip}{0pt}
\vspace*{8mm}
{\sffamily\bfseries\small\color{rubricCrimson}REPLACE BEFORE SUBMISSION\par}
\vspace{4mm}
{\fontsize{26}{31}\selectfont\bfseries Sample marking rubric\par}
\vspace{4mm}
These criteria and point values illustrate how a marking sheet can be laid out.
Use the rubric supplied for your assignment when preparing the final report.\par

\vspace{13mm}
\noindent\begin{tblr}{
  width=\linewidth,
  colspec={X[1.15,l] X[1.6,l] Q[c,wd=1.55cm] Q[c,wd=1.75cm]},
  column{1}={leftsep=0pt}, column{4}={rightsep=0pt},
  rows={m,rowsep=6mm}, row{1,Z}={font=\bfseries,rowsep=3mm},
  hline{1,Z}={0.8pt}, hline{2,3,4,5}={0.35pt},
}
  Criterion & Sample evidence & Possible & Awarded \\
  Criterion A & Example outline & 4 & \\
  Criterion B & Example explanation & 4 & \\
  Criterion C & Example illustration & 2 & \\
  Total & & 10 & \\
\end{tblr}\par

\vspace{16mm}
{\bfseries Marker comments}\par
\vspace{12mm}\hrule height 0.35pt
\vspace{12mm}\hrule height 0.35pt
\vspace{12mm}\hrule height 0.35pt

\vfill
{\small To replace this page, save the course PDF as
\texttt{local-materials/rubric.pdf} and recompile.\par}
\clearpage
\endgroup
}

\end{document}
